\section*{章节综合：预测损失也是一种编码代价}
假设样本中第一类占四分之一，用二元模型概率 $q$ 来编码或预测。经验平均交叉熵为
\[C(q)=-\tfrac14\log_2q-\tfrac34\log_2(1-q),\qquad 0<q<1.\]
求导得到 $C'(q)=(q-\frac14)/(q(1-q)\ln2)$。分母为正，所以 $q<1/4$ 时下降，$q>1/4$ 时上升，最小点恰在经验频率处。

这说明在当前二元、独立样本模型中，匹配频率会最小化经验编码代价。它没有保证未来样本仍有相同频率，也没有处理参数受复杂网络约束的情况。经验最优与总体最优的区别，需要统计推断补足。

在实验室把 $p$ 固定为 $0.25$，拖动 $q$，核对导数符号所预测的升降与最低点。再检查最低值等于真实熵，而多付的部分等于 KL。通过概率、编码长度、优化导数三种语言理解同一问题，便完成了这两课的连接。
